%%%%%%%%%%%%%%%%%%%%%%%%%%%%%%%%%%%%%%%%%%%%%%%%%%%%%%%%%%%%%%%%%%%%%%%%%%%%%%%%%%%%%%%%%%%%%%%%%%%%%%%%%%%%%%%%%%%%%%%%%%%%%%%%%%%%%%%%%%%%%%%%%%%%%%%%%%%
% This is just an example/guide for you to refer to when producing your supplementary material for your Frontiers article.                                 %
%%%%%%%%%%%%%%%%%%%%%%%%%%%%%%%%%%%%%%%%%%%%%%%%%%%%%%%%%%%%%%%%%%%%%%%%%%%%%%%%%%%%%%%%%%%%%%%%%%%%%%%%%%%%%%%%%%%%%%%%%%%%%%%%%%%%%%%%%%%%%%%%%%%%%%%%%%%

%%% Version 2.5 Generated 2018/06/15 %%%
%%% You will need to have the following packages installed: datetime, fmtcount, etoolbox, fcprefix, which are normally inlcuded in WinEdt. %%%
%%% In http://www.ctan.org/ you can find the packages and how to install them, if necessary. %%%
%%%  NB logo1.jpg is required in the path in order to correctly compile front page header %%%

\documentclass[utf8]{frontiers_suppmat} % for all articles
\usepackage{amssymb}
\usepackage{amsmath, bm}
\usepackage{siunitx}
\usepackage{url,hyperref,lineno,microtype}
\let\DeclareUSUnit\DeclareSIUnit
\let\US\SI
\DeclareUSUnit\inch{in}
\providecommand{\abs}[1]{\lvert#1\rvert}
\providecommand{\norm}[1]{\lVert#1\rVert}
\usepackage{graphicx}
\graphicspath{{../}}
\usepackage{caption}
\usepackage{tabularx}
\usepackage{multirow}
\usepackage{array}
\usepackage[onehalfspacing]{setspace}
\renewcommand{\arraystretch}{1.5}  % Increase row height (default is 1.0)
%\linenumbers

% Leave a blank line between paragraphs instead of using \\

\begin{document}
\onecolumn
\firstpage{1}
\title[Supplementary Material]{{\helveticaitalic{Supplementary Material}}}
\maketitle

\section{Supplementary Tables and Figures}

\subsection{Tables}

\begin{table}[htbp]
\caption{Comparison of pre- and post-optimization torque prediction fits for BPAs in flexor and extensor configurations on the pinned and biomimetic knees. GoF metrics are RMSE ($N \cdot m$), FVU, and maximum residual ($N \cdot m$).}
\label{tab:TorqueFit}
\resizebox{\textwidth}{!}
{%
\begin{tabular}{|cc|l|l|lll|lll|}
\hline
\multicolumn{2}{|l|}{\multirow{2}{*}{\textbf{Configuration}}} &
  \multirow{2}{*}{\textbf{$\phi$}} &
  \multirow{2}{*}{\textbf{$l_{rest}$, cm}} &
  \multicolumn{3}{c|}{\textbf{Pre-optimization}} &
  \multicolumn{3}{c|}{\textbf{Optimized}} \\ \cline{5-10} 
\multicolumn{2}{|l|}{} &
   &
   &
  \multicolumn{1}{l|}{\textbf{RMSE}} &
  \multicolumn{1}{l|}{\textbf{FVU}} &
  \textbf{Max. Residual} &
  \multicolumn{1}{l|}{\textbf{RMSE}} &
  \multicolumn{1}{l|}{\textbf{FVU}} &
  \textbf{Max. Residual} \\ \hline
\multicolumn{1}{|c|}{\multirow{4}{*}{\textbf{Flexor}}} &
  \multirow{2}{*}{\textbf{Pinned}} &
  10 &
  48.5 &
  \multicolumn{1}{l|}{4.871} &
  \multicolumn{1}{l|}{0.723} &
  9.865 &
  \multicolumn{1}{l|}{1.608} &
  \multicolumn{1}{l|}{0.095} &
  3.242 \\ \cline{3-10} 
\multicolumn{1}{|c|}{} &
   &
  10 &
  45.7 &
  \multicolumn{1}{l|}{3.024} &
  \multicolumn{1}{l|}{0.163} &
  6.753 &
  \multicolumn{1}{l|}{1.915} &
  \multicolumn{1}{l|}{0.072} &
  4.695 \\ \cline{2-10} 
\multicolumn{1}{|c|}{} &
  \multirow{2}{*}{\textbf{Biomimetic}} &
  10 &
  41.5 &
  \multicolumn{1}{l|}{7.448} &
  \multicolumn{1}{l|}{5.557} &
  13.308 &
  \multicolumn{1}{l|}{0.680} &
  \multicolumn{1}{l|}{0.052} &
  1.055 \\ \cline{3-10} 
\multicolumn{1}{|c|}{} &
   &
  20 &
  41.5 &
  \multicolumn{1}{l|}{14.803} &
  \multicolumn{1}{l|}{3.366} &
  23.715 &
  \multicolumn{1}{l|}{2.868} &
  \multicolumn{1}{l|}{0.126} &
  3.550 \\ \hline
\multicolumn{1}{|c|}{\multirow{4}{*}{\textbf{Extensor}}} &
  \multirow{3}{*}{\textbf{Pinned}} &
  10 &
  41.5 &
  \multicolumn{1}{l|}{4.520} &
  \multicolumn{1}{l|}{3.972} &
  11.741 &
  \multicolumn{1}{l|}{2.042} &
  \multicolumn{1}{l|}{0.844} &
  3.588 \\ \cline{3-10} 
\multicolumn{1}{|c|}{} &
   &
  10 &
  45.7 &
  \multicolumn{1}{l|}{1.999} &
  \multicolumn{1}{l|}{0.892} &
  4.609 &
  \multicolumn{1}{l|}{0.896} &
  \multicolumn{1}{l|}{0.179} &
  1.846 \\ \cline{3-10} 
\multicolumn{1}{|c|}{} &
   &
  10 &
  48 &
  \multicolumn{1}{l|}{0.845} &
  \multicolumn{1}{l|}{0.232} &
  1.951 &
  \multicolumn{1}{l|}{0.682} &
  \multicolumn{1}{l|}{0.151} &
  2.068 \\ \cline{2-10} 
\multicolumn{1}{|c|}{} &
  \multicolumn{1}{l|}{\textbf{Biomimetic}} &
  10 &
  51.8 &
  \multicolumn{1}{l|}{3.032} &
  \multicolumn{1}{l|}{4.486} &
  5.587 &
  \multicolumn{1}{l|}{1.461} &
  \multicolumn{1}{l|}{1.041} &
  2.755 \\ \hline
\end{tabular}
}
\end{table}

\begin{table}[htbp]
\caption{Muscle origin and insertion locations for different models. Origin is in femur reference frame, insertion is in the tibia reference frame.}
\label{table:muscle_attachment}
\centering
\begin{tabular}{||l l l r r r||} 
 \hline
 Muscle & Origin/Insertion & Model & X (m) & Y (m) & Z (m)  \\ [0.5ex] 
 \hline\hline
 Flexor & origin & Pinned knee & -0.075 & 0.100 & 0.0328 \\
 Flexor & insertion & Pinned knee & -0.05011 & -0.045 & 0.0326 \\
 Extensor & origin & Pinned knee & 0.030 & -0.050 & 0 \\
 Extensor & insertion & Pinned knee & 0.0425 & -0.076 & 0 \\  
 Bifemsh & origin & Gait2392 & 0.005 & -0.211 & 0.023 \\
 Bifemsh & insertion & Gait2392 & -0.03 & -0.036 & 0.029 \\
 Vastus int. & origin & Gait2392 & 0.029 & -0.192 & 0.031 \\
 Vastus int. & insertion & Gait2392 & moving & moving & 0.0018 \\  
 Flexor, \qty{10}{\milli\metre} & origin & Biomimetic knee & -0.050 & -0.045 & 0.0328 \\
 Flexor, \qty{10}{\milli\metre} & insertion & Biomimetic knee  & -0.0279 & -0.046 & 0.0328 \\
 Extensor, \qty{10}{\milli\metre} & origin & Biomimetic knee  & 0.040 & 0.035 & 0 \\
 Extensor, \qty{10}{\milli\metre} & insertion & Biomimetic knee  & 0.021 & -0.072 & 0 \\
 Flexor, \qty{20}{\milli\metre} & origin & Biomimetic knee  & -0.050 & -0.045 & 0.0328  \\
 Flexor, \qty{20}{\milli\metre} & insertion & Biomimetic knee  & -0.0279 & -0.046 & 0.0328 \\
 \hline
\end{tabular}
\end{table}

\clearpage

\subsection{BPA routing and via points}

The tables in this section list the routing via points of the BPA configurations reported in the main text. All coordinates are coded in meters in the project's MATLAB sources and are reported here in millimeters, rounded to $0.01$~mm. Knee angle $\theta_{k}$ follows the model convention with flexion negative. Femur-side points (path rows before the crossing index) are specified in the femur-referenced instantaneous-center-of-rotation (ICR) frame and are used untransformed at every knee angle, while tibia-side points are specified in the tibia $\theta_{1}$ frame and are mapped into the ICR frame at each pose with the measured ICR migration, fitted as cubic splines of the tibia-to-ICR offsets over the knee range. The pinned-knee bracket reference points use the knee-ICR and hip frames noted per row, and the pinned-knee route endpoints (origins and insertions) are listed in Table~\ref{table:muscle_attachment}.

\begin{table}[htbp]
\caption{Bracket reference points of the pinned-knee BPA routes, in mm. The pinned flexor route runs from origin to insertion with a compliant bracket at each end, and the pinned extensor route additionally wraps on two fixed arcs ($R_{1} = \qty{12}{\milli\metre}$, active for $-23^{\circ} \leq \theta_{k} \leq 27^{\circ}$, and $R_{2} = \qty{40}{\milli\metre}$, active for $\theta_{k} \leq -23^{\circ}$).}
\label{tab:RoutePinnedBrackets}
\centering
\small
\begin{tabular}{lllrrr}
\hline
\textbf{Configuration} & \textbf{Point} & \textbf{Frame} & \textbf{x (mm)} & \textbf{y (mm)} & \textbf{z (mm)} \\
\hline
Pinned flexor & $\mathbf{p}_{bri}$, insertion bracket & knee ICR & -48.11 & -107.81 & 13.80 \\
Pinned flexor & $\mathbf{p}_{br2}$, origin-side second bracket & hip & -52.61 & 0.00 & 75.06 \\
Pinned extensor & $\mathbf{p}_{br}$, rib midpoint, medial side & hip & -3.84 & -46.44 & 62.50 \\
\hline
\end{tabular}
\end{table}

\begin{table}[htbp]
\caption{Routing via points of the biomimetic-knee \qty{10}{\milli\metre} BPA routes, in mm. The flexor route is a two-point path (crossing index 2), and the extensor route is the eight-point path of the extensor test (crossing index 5, nominal resting length \qty{520}{\milli\metre}) whose femur- and tibia-side contacts engage progressively with flexion. The active-range column lists the knee angles over which each point participates in the path.}
\label{tab:RouteBio10}
\centering
\small
\begin{tabular}{lllrrrl}
\hline
\textbf{Point} & \textbf{Description} & \textbf{Frame} & \textbf{x (mm)} & \textbf{y (mm)} & \textbf{z (mm)} & \textbf{Active $\theta_{k}$} \\
\hline
\multicolumn{7}{l}{\textbf{Flexor, \qty{10}{\milli\metre} (two-point route, resting length \qty{415}{\milli\metre})}} \\
\hspace{3mm} $p_{1}$ & origin & femur & -50.00 & 35.00 & 50.00 & all $\theta_{k}$ \\
\hspace{3mm} $p_{2}$ & insertion & tibia $\theta_{1}$ & -12.24 & -8.87 & 27.87 & all $\theta_{k}$ \\
\multicolumn{7}{l}{\textbf{Extensor, \qty{10}{\milli\metre} (eight-point route, resting length \qty{520}{\milli\metre})}} \\
\hspace{3mm} $p_{1}$ & origin & femur & 40.00 & 35.00 & 0.00 & all $\theta_{k}$ \\
\hspace{3mm} $p_{2}$ & mounting-base contact & femur & 75.90 & -274.76 & 0.00 & all $\theta_{k}$ \\
\hspace{3mm} $p_{3}$ & femur channel contact & femur & 59.82 & -374.27 & 0.00 & $\theta_{k} \leq -9.19^{\circ}$ \\
\hspace{3mm} $p_{4}$ & femoral condyle contact & femur & 39.55 & -421.83 & 0.00 & $\theta_{k} \leq -69.8^{\circ}$ \\
\hspace{3mm} $p_{5}$ & tibia contact & tibia $\theta_{1}$ & 58.71 & 25.00 & 0.00 & $\theta_{k} \leq -9.19^{\circ}$ \\
\hspace{3mm} $p_{6}$ & tibia contact & tibia $\theta_{1}$ & 58.71 & 12.28 & 0.00 & all $\theta_{k}$ \\
\hspace{3mm} $p_{7}$ & tibia tendon contact & tibia $\theta_{1}$ & 54.00 & -6.12 & 0.00 & all $\theta_{k}$ \\
\hspace{3mm} $p_{8}$ & patellar ligament ring & tibia $\theta_{1}$ & 36.04 & -28.44 & 0.00 & all $\theta_{k}$ \\
\hline
\end{tabular}
\end{table}

\begin{table}[htbp]
\caption{Routing via points of the \qty{20}{\milli\metre} route redesigns, in mm. The flexor design routes two BPAs on mirrored paths: the second BPA keeps the $z$-mirrored distal endpoint while both origins sit off the same side of the knee, and each route carries one tibial wrap point placed at the BPA centerline standoff from the tibial surface while contact is active. The extensor design is the clip-on-cam route; the liftoff column lists the knee angle above which each contact leaves the path as the knee extends, with all contacts engaged at maximum flexion.}
\label{tab:RouteRedesign20}
\centering
\small
\begin{tabular}{lllrrrl}
\hline
\textbf{Point} & \textbf{Description} & \textbf{Frame} & \textbf{x (mm)} & \textbf{y (mm)} & \textbf{z (mm)} & \textbf{Liftoff $\theta_{k}$} \\
\hline
\multicolumn{7}{l}{\textbf{Flexor, \qty{20}{\milli\metre} (two-BPA mirrored route, two BPAs, resting length \qty{561.49}{\milli\metre}, tendon \qty{25.00}{\milli\metre})}} \\
\hspace{3mm} $p_{1}$ & origin, BPA 1 & femur & -125.00 & 135.00 & 72.08 & end point \\
\hspace{3mm} $p_{end}$ & insertion, BPA 1 & tibia $\theta_{1}$ & -18.72 & -32.69 & 38.96 & end point \\
\hspace{3mm} $p_{1}^{(2)}$ & origin, BPA 2 & femur & -125.00 & 135.00 & -5.85 & end point \\
\hspace{3mm} $p_{end}^{(2)}$ & insertion, BPA 2 & tibia $\theta_{1}$ & -18.72 & -32.69 & -38.96 & end point \\
\multicolumn{7}{l}{\textbf{Extensor, \qty{20}{\milli\metre} (clip-on-cam route, two BPAs, resting length \qty{827.74}{\milli\metre}, tendon \qty{92.49}{\milli\metre})}} \\
\hspace{3mm} $p_{1}$ & origin (fixed) & femur & 40.00 & 137.00 & 0.00 & end point \\
\hspace{3mm} $p_{2}$ & mounting-base contact & femur & 83.00 & -271.70 & 0.00 & $-21.52^{\circ}$ \\
\hspace{3mm} $p_{3}$ & femur contact & femur & 66.60 & -427.80 & 0.00 & $-54.34^{\circ}$ \\
\hspace{3mm} $p_{4}$ & femur contact & femur & 54.70 & -443.60 & 0.00 & $-89.80^{\circ}$ \\
\hspace{3mm} $p_{5}$ & femoral condyle contact & femur & 38.80 & -454.40 & 0.00 & $-110.81^{\circ}$ \\
\hspace{3mm} $p_{6}$ & cam contact (optimized) & tibia $\theta_{1}$ & 116.32 & 15.04 & 0.00 & $-18.89^{\circ}$ \\
\hspace{3mm} $p_{7}$ & cam contact (optimized) & tibia $\theta_{1}$ & 131.32 & -18.97 & 0.00 & none \\
\hspace{3mm} $p_{8}$ & cam contact (optimized) & tibia $\theta_{1}$ & 127.04 & -47.87 & 0.00 & none \\
\hspace{3mm} $p_{9}$ & distal ring (insertion) & tibia $\theta_{1}$ & 34.42 & -197.10 & 0.00 & end point \\
\hline
\end{tabular}
\end{table}

\clearpage

\subsection{Figures}

\begin{figure}[htbp]
\begin{center}
\includegraphics[width=1\textwidth]{Figures/RouteVariants_20261020.eps}
\end{center}
\caption{Predicted knee-axis torque of the redesigned route families against the human-scaled requirements. Knee angle is negative in flexion; extension torque is positive and flexion torque negative. \textbf{(A)} Extensor routes: the direct route and the cam design of record against the vasti requirement and its required \qty{5}{\percent} envelope. \textbf{(B)} Flexor routes: the direct two-BPA route and the mono-pulley variant on the insertion-side measure against the biceps femoris short-head requirement and its envelope.}\label{fig:RouteVariants}
\end{figure}

\begin{figure}[htbp]
\begin{center}
\includegraphics[width=1\textwidth]{Figures/extensor_misalign.png}
\end{center}
\caption{During extensor tests with the pinned knee using a tendon it was noticed that the actuator pathway was modeled to be anterior of the socket head cap screw that joins the femur body to the knee assembly. This slippage affects the muscle path and the moment arm. It was determined this caused a \qty{\pm20}{\milli\metre} $\hat{z}$ axis displacement. Test assembly is shown with \textbf{(A)} top down view and  \textbf{(B)} side view.}
\label{fig:extensor_misalign}
\end{figure}

\bibliography{BolenFrontiers22}

\end{document}
